\section{Cronograma de actividades}

El semestre inicia el 14 de julio de 2026 (semana 1) y la sustentación se realiza a más tardar el 20 de noviembre de 2026 (semana 18).

\begin{table}[H]
\centering
\singlespacing
\scriptsize
\caption{Cronograma de actividades por quincena (período académico 2026-2).}
\label{tab:cronograma}
\setlength{\tabcolsep}{3pt}
\renewcommand{\arraystretch}{0.95}
\begin{tabular}{p{5.2cm}*{9}{c}}
\toprule
\textbf{Actividad} & \textbf{1--2} & \textbf{3--4} & \textbf{5--6} & \textbf{7--8} & \textbf{9--10} & \textbf{11--12} & \textbf{13--14} & \textbf{15--16} & \textbf{17--18}\\
\midrule
Revisión bibliográfica y anteproyecto & X & X & & & & & & & \\
F1. Caracterización térmica y \emph{springback} & & X & X & X & & & & & \\
F2. Diseño del subsistema térmico & & & X & X & & & & & \\
F3. Diseño mecánico y CAD & & & & X & X & & & & \\
F4. Control PID e interfaz táctil & & & & & X & X & X & & \\
F5. Fabricación, integración y montaje & & & & & & X & X & & \\
F6. Validación experimental & & & & & & & X & X & \\
Informes de avance y final, y entrega al laboratorio & & & & & X & & X & X & \\
Preparación y sustentación pública & & & & & & & & & X\\
\bottomrule
\end{tabular}
\end{table}
